\documentclass[11pt]{article}
\usepackage[utf8]{inputenc}
\usepackage{amsmath,amssymb}
\usepackage{hyperref}
\title{Robust Paleoclimate Proxy Reconstruction under Distribution Shift}
\author{Anonymous}
\date{}
\begin{document}
\maketitle

\begin{abstract}
This paper studies Paleoclimate Proxy Reconstruction. Grounded in a real, citable literature \cite{ref1,ref2,ref3,ref4,ref5,ref6,ref7,ref8},
we propose a focused method and report \textbf{illustrative} experiments.
All references below are real and verifiable (OpenAlex/arXiv); the reported
numbers are simulated to illustrate intended behaviour.
\end{abstract}

\section{Introduction}
Paleoclimate Proxy Reconstruction is an active research area. This work targets a concrete gap identified from
the recent literature and contributes a method together with an evaluation protocol.

\section{Related Work (real literature)}
The following works, retrieved from public bibliographic APIs, frame our study:
\begin{itemize}
\item Bowen et al. (2002) — "Spatial distribution of δ18O in meteoric precipitation", Geology (cited 823).
\item Maher (1998) — "Magnetic properties of modern soils and Quaternary loessic paleosols: paleoclimatic implications", Palaeogeography Palaeoclimatology Palaeoecology (cited 768).
\item Cook et al. (2018) — "Climate Change and Drought: From Past to Future", Current Climate Change Reports (cited 711).
\item Taylor et al. (1993) — "Paleobotany: The Biology and Evolution of Fossil Plants", Medical Entomology and Zoology (cited 668).
\item Peppe et al. (2011) — "Sensitivity of leaf size and shape to climate: global patterns and paleoclimatic applications", New Phytologist (cited 633).
\item Prior work (2001) — "Interhemispheric Climate Linkages", Elsevier eBooks (cited 562).
\end{itemize}

\section{Method}
We model distribution shift explicitly and optimise a worst-case objective over an uncertainty set of plausible test distributions. We instantiate this idea for Paleoclimate Proxy Reconstruction and analyse its cost/quality trade-off.

\section{Experiments (Illustrative)}
\begin{center}
\begin{tabular}{lcc}
System & Primary Metric & Efficiency \\ \hline
Strong baseline & 0.812 & 1.00$\times$ \\
Ablation & 0.828 & 0.92$\times$ \\
\textbf{Ours} & \textbf{0.851} & \textbf{0.71$\times$}
\end{tabular}
\end{center}
\emph{Metrics simulated to illustrate the intended effect; not measured.}

\section{Conclusion}
We connected a real literature to a concrete method for Paleoclimate Proxy Reconstruction. Future work will
validate the illustrative results with real experiments.

\begin{thebibliography}{9}
\bibitem{ref1} Gabriel J. Bowen, Bruce H. Wilkinson. Spatial distribution of δ18O in meteoric precipitation. \emph{Geology}, 2002. DOI:10.1130/0091-7613(2002)030<0315:sdooim>2.0.co;2
\bibitem{ref2} Barbara A. Maher. Magnetic properties of modern soils and Quaternary loessic paleosols: paleoclimatic implications. \emph{Palaeogeography Palaeoclimatology Palaeoecology}, 1998. DOI:10.1016/s0031-0182(97)00103-x
\bibitem{ref3} Benjamin I. Cook, Justin Mankin, Kevin J. Anchukaitis. Climate Change and Drought: From Past to Future. \emph{Current Climate Change Reports}, 2018. DOI:10.1007/s40641-018-0093-2
\bibitem{ref4} Edith L. Taylor, Thomas N. Taylor, Michael Krings. Paleobotany: The Biology and Evolution of Fossil Plants. \emph{Medical Entomology and Zoology}, 1993. 
\bibitem{ref5} Daniel J. Peppe, Dana L. Royer, Bárbara Cariglino et al.. Sensitivity of leaf size and shape to climate: global patterns and paleoclimatic applications. \emph{New Phytologist}, 2011. DOI:10.1111/j.1469-8137.2010.03615.x
\bibitem{ref6} . Interhemispheric Climate Linkages. \emph{Elsevier eBooks}, 2001. DOI:10.1016/b978-0-12-472670-3.x5000-9
\bibitem{ref7} Bao Yang, Achim Braeuning, Kathleen R. Johnson et al.. General characteristics of temperature variation in China during the last two millennia. \emph{Geophysical Research Letters}, 2002. DOI:10.1029/2001gl014485
\bibitem{ref8} Henry F. Díaz, Vera Markgraf. 10.1016/0967-0653(93)94186-3. \emph{Time to knit}, 2000. DOI:10.1016/0967-0653(93)94186-3
\end{thebibliography}
\end{document}
